% Generated by tools/build_v6.py; edit the assembly rules, not this file.
\section{A uniform variance--work upper constant of 713}\label{sec:713}
The finite-budget construction in \cref{sec:classifier,sec:minimax-proof}
uses the original threshold and constants $85,149$. For the unbiased
upper bound in \cref{thm:unbiased}, change the classifier threshold to
\[
 t_{n,N}=\sqrt{\frac{1/4+\log(N/n)}n}\qquad(n<N),
\]
keeping the same dyadic checkpoints, strict crossing rule, negative final
tie convention, shared-prefix increments, and independent level law
$p_\ell=(1-r)r^\ell$, $r=2^{-1/2}$. In this section $Z_N,\tau_N,
\Delta_\ell,Y_L,\mathcal W_L$ denote this modified construction only.
It observes no magnitude or sign and uses no distributional parameters.

\begin{theorem}[Separate conditional moment bounds]\label{thm:713}
For every $0<d\le1/2$, either sign, and finite $L$, $N=2^L$,
\[
 \E[Y_L^2\mid d,S]\le23\min\{\sqrt N,d^{-1}\},\qquad
 \E[\mathcal W_L\mid d,S]\le31\min\{\sqrt N,d^{-1}\}.
\]
If $I_\mu<\infty$, the untruncated estimator is square-integrable,
unbiased, and has finite expected actual work, with
\[
 \sup_q\E_{\mu,q}Y^2\le23I_\mu,\qquad
 \sup_q\E_{\mu,q}\mathcal W\le31I_\mu.
\]
Consequently the product of the separate variance and work suprema is
at most $713I_\mu^2$.
\end{theorem}

\paragraph{Error envelope.}
An incorrect early stop at prefix $n$ requires a one-sided deviation of
at least $d+t_{n,N}$. Hoeffding's inequality and a union bound over the
fixed dyadic prefixes yield, with $x=Nd^2$,
\[
 e_N\le e^{-2x}+e^{-1/2}N^{-2}
             \sum_{n<N,\ n\text{ dyadic}}n^2e^{-2nd^2}
 \le e^{-2x}+\tfrac23e^{-1/2}x^{-2}.
\]
For the second step integrate $u e^{-bu}$ over the disjoint intervals
$[n/2,n]$: each integral is at least $3n^2e^{-bn}/8$, so the sum is
at most $8/(3b^2)$. Since $x^2e^{-2x}\le e^{-2}$ for every $x>0$,
\[
 e_N\le\min\{1,A x^{-2}\},\qquad
 A=e^{-2}+\tfrac23e^{-1/2}<\tfrac17+\tfrac5{12}
 =\tfrac{47}{84}<\tfrac35.
\]
Here $e^2>7$ and $e^{1/2}>8/5$ follow from their positive power series.
For $\ell\ge1$, a nonzero increment requires a fine or coarse error;
the coarse scale has $x/2$. Hence
\begin{equation}\label{eq:713-increment}
 \E\Delta_\ell^2\le\min\{1,5A x^{-2}\}
             \le\min\{1,3x^{-2}\}.
\end{equation}
The same envelope covers $\ell=0$, because $d^2\le1/4$ and
$\Delta_0=Z_1\in\{0,1\}$.

\paragraph{Second moment.}
Put $d'=d/3^{1/4}$. Summing \eqref{eq:713-increment} divided by
$p_\ell$ and splitting at $2^\ell(d')^2=1$ gives
\[
 \E Y_L^2\le C_0\min\{\sqrt{2^L},1/d'\},\qquad
 C_0=\frac1{1-r}\left(\frac1{1-r}+\frac1{1-r^3}\right)<17.
\]
Before the split the terms are geometric with size $\sqrt{2^\ell}$;
after it their ratio is $r^3$. This includes cutoffs before the split.
The rational comparison $r<442/625$ and monotonicity of the displayed
expression prove $C_0<17$ by integer arithmetic. Since
$3^{1/4}<4/3$, the conditional constant is less than $68/3<23$.

\paragraph{Expected work.}
Take $c=13/7$. Survival at a dyadic prefix $n<N$ and $d>ct_{n,N}$
requires a one-sided deviation larger than $(1-1/c)d$. Therefore
\[
 \Prob(\tau_N>n)\le\ind\{d\le ct_{n,N}\}
            +\exp\{-2(1-1/c)^2nd^2\}.
\]
Suppose $x=Nd^2\ge1$ and write $u=nd^2$. The indicator implies
$u\le c^2[1/4+\log(x/u)]$. If $u<1$, then
$u\le1+c^2\log x$; if $u\ge1$, use $\log u\ge0$ and
$c^2/4=169/196<1$ for the same conclusion. Its dyadic contribution
is at most $2d^{-2}(1+c^2\log x)$.
For $b>0$, integration on $[n/2,n]$ gives
$\sum n e^{-bnd^2}\le2/(bd^2)$. The survival identity and
$1\le d^{-2}/4$ now give
\[
 \E\tau_N
 \le d^{-2}\left[\frac{125}{18}+\frac{338}{49}\log x\right]
 \le7d^{-2}(1+\log x).
\]
Together with $\tau_N\le N$, this proves $\E\tau_N\le7m_N(d)$.

On a shared stream the cap-$N$ threshold is at least the cap-$N/2$
threshold at every prefix before $N/2$. If the fine classifier stops
there, the coarse one has already stopped; at or after $N/2$ the coarse
one has reached its cap. Thus $\tau_{N/2}\le\tau_N$ pathwise and the
actual work of the coupled increment equals $\tau_N$. Level zero
costs one. It follows that
\[
 \E\mathcal W_L\le7(1-r)\sum_{\ell=0}^L
          2^{-\ell/2}m_{2^\ell}(d)
 \le7\left(2+\frac{\log2}{1-r}\right)
          \min\{\sqrt{2^L},d^{-1}\}.
\]
To obtain the last inequality, split at $2^\ell d^2=1$ as in
\cref{lem:universal}. The first tail value satisfies $1<2^\ell d^2\le2$;
the remaining sum is bounded by
$d^{-1}\sum_{k\ge0}r^k(1+\log2+k\log2)$.
Using $r<71/100$ and $\log2<7/10$ bounds the final constant by
$896/29<31$.

\paragraph{Untruncated limit and separate suprema.}
Integrate the conditional bounds over $(D,S)$. Monotone convergence of
the nonnegative moment and cost series gives finite second moment and
work when $I_\mu<\infty$. More explicitly,
\[
 \E[(Y-Y_L)^2]=\sum_{\ell>L}\frac{\E\Delta_\ell^2}{p_\ell}
       \longrightarrow0.
\]
Thus $Y_L\to Y$ in $L^2$ and $L^1$. Finite telescoping gives
$\E Y_L=\E Z_{2^L}$; the conditional error tends to zero for every
$d>0$, and bounded convergence proves $\E Y=\theta$.
The bounds hold for every $q$ with the same estimator and the same
$I_\mu$. Taking the two suprema separately and using
$\Var Y\le\E Y^2$ proves \cref{thm:713} and the improved upper
constant in \cref{thm:unbiased}. No interchange of a supremum with a
strict inequality is needed.

\paragraph{Formal scope.}
\observableleanref{observable\_713} realizes one geometric-seed estimator
in the original observable adaptive experiment, before every gap law.
Its separate extended variance and actual work bounds are $23I_\mu$ and
$31I_\mu$, uniformly over every measurable sign law. The product of the
two separate suprema is at most $713I_\mu^2$. Both classifiers read the
same coin's observed prefix; exact execution identities transfer their
output and flip count to the infinite experiment. Together with the
converse, \observableleanref{observable\_unbiased\_iff} proves the full
finite-variance/work existence characterization in this model.

The earlier, more general
\leanref{Uniform713/FullBound.lean}{full\_bound\_713} proves the
untruncated estimator result in a conditional observation-kernel model,
including bounded independent observations with means $1/2\pm D$.
\leanref{ResearchRouteA/SeparateBounds.lean}{uniform\_family\_bounds}
supplies the separate-supremum conclusion for a fixed kernel and
estimator over a nonempty family of latent laws sharing $I_\mu$.
The finite-budget calibration wrapper for the original RMS law has its
separate realization in \cref{sec:verification}.
